\begin{definition}[Ordered coefficient of a shared-side insertion]%
    \label{def:peps_torus_shared_flux_ordered_element}
    \lean{TNLean.PEPS.torusSharedFluxOrderedElement}
    \leanok
    For a plaquette based at $v=(x,y)$, write $k_v(g)=g$ if the shared
    upward edge from $(x+1,y)$ to $(x+1,y+1)$ agrees with the native
    endpoint order, and $k_v(g)=g^{-1}$ otherwise. Thus its directed
    upward transport is always $g$, including the vertical seam.
\end{definition}

\begin{theorem}[Literal opposite fluxes on adjacent plaquettes]%
    \label{thm:peps_torus_shared_flux_pair}
    \lean{TNLean.PEPS.regularWalkHolonomy_torusSharedPlaquetteInsertion,
        TNLean.PEPS.torusSharedPlaquetteInsertion_eq_treeCycleAssignment}
    \leanok
    \uses{def:peps_torus_shared_flux_ordered_element,
        thm:peps_translated_two_plaquette_geometry}
    On a finite torus of horizontal period at least four and vertical period
    at least three, insert directed transport $g$ on the shared upward side
    of the two plaquettes based at $(x,y)$ and $(x+1,y)$, and the identity
    on every other bond. Their actual counterclockwise holonomies are
    $(g,g^{-1})$. The ordered insertion is precisely $k_v(g)$ on the
    derived middle chord of the translated six-site spanning tree.
    These assertions hold at every position, including both seams.
\end{theorem}
\begin{proof}\leanok
    The shared upward side of the left plaquette is the reversed left side
    of the right plaquette. The remaining transports are identities.
    Compare the ordered coefficient with the translated middle chord.
\end{proof}

\begin{theorem}[Coherent opposite-pair creation on six original spins]%
    \label{thm:peps_torus_translated_flux_creation}
    \lean{TNLean.PEPS.IsGIsometric.exists_unitary_torusTranslatedPhysicalFluxCreation,
        TNLean.PEPS.IsGIsometric.exists_unitary_torusTranslatedGlobalFluxCreation}
    \leanok
    \uses{thm:peps_torus_shared_flux_pair,
        thm:peps_regular_physical_flux_creation,
        thm:peps_regular_cycle_global_flux_creation}
    Let the original four-leg tensor be $G$-isometric for the regular
    action. For every position $v$ and class representative $g$, there is
    a unitary $W_{v,g}$ on the six original spins surrounding these two
    plaquettes such that, for every outgoing boundary column $\theta$,
    \begin{align}
        W_{v,g}\Psi_R(1,\theta)
        =\gamma_g\sum_{z\in G}\Psi_R(zgz^{-1},\theta),
        \qquad \gamma_g=(|G|\,|C_G(g)|)^{-1/2}.
    \end{align}
    In each summand the shared upward side has transport $zgz^{-1}$ and
    every other internal bond has identity transport. There is also a six-spin
    unitary whose complementary identity extension is globally unitary, with
    the same coherent-sum action on the actual globally contracted states,
    for every common exterior and crossing assignment. The unitary is
    chosen before all boundary and exterior data.

    The literal summands have adjacent flux pairs
    $(zgz^{-1},zg^{-1}z^{-1})$. This is the normalized regular finite-torus
    creation operation of \cite[Theorem~6.17]{Schuch2010PEPS} for periods
    at least four and three. Its parent-Hamiltonian interpretation and
    unrestricted geometry remain separate; see
    \cite{gap:rmp_peps_quantum_double_g_isometry}.
\end{theorem}
\begin{proof}\leanok
    Apply the actual chosen-cycle creation theorem to the derived middle
    chord, using the ordered element $k_v(g)$. Inversion commutes with
    conjugation, so its coherent insertions are the directed insertions
    displayed above. The centralizers of $g$ and $g^{-1}$ coincide; hence
    the seam does not change $\gamma_g$. The actual cut transport proves
    the global statement with the same common exterior operators.
\end{proof}
